\documentclass[12pt,a4paper]{report}

% ── Encodage & langue ──────────────────────────────────────────────────────────
\usepackage[utf8]{inputenc}
\usepackage[T1]{fontenc}
\usepackage[french]{babel}

% ── Mise en page ───────────────────────────────────────────────────────────────
\usepackage[top=2.5cm, bottom=2.5cm, left=3cm, right=2.5cm]{geometry}
\usepackage{setspace}
\onehalfspacing

% ── Graphiques & couleurs ──────────────────────────────────────────────────────
\usepackage{graphicx}
\usepackage{xcolor}
\definecolor{ensetBlue}{RGB}{0,70,127}
\definecolor{lightGray}{RGB}{245,245,245}
\definecolor{commentGreen}{RGB}{80,161,79}
\definecolor{keywordBlue}{RGB}{86,156,214}
\definecolor{stringOrange}{RGB}{206,145,120}

% ── Hyperliens ─────────────────────────────────────────────────────────────────
\usepackage[colorlinks=true, linkcolor=ensetBlue, urlcolor=ensetBlue]{hyperref}

% ── Listings (code source) ─────────────────────────────────────────────────────
\usepackage{listings}
\lstdefinestyle{javaStyle}{
    backgroundcolor=\color{lightGray},
    commentstyle=\color{commentGreen}\itshape,
    keywordstyle=\color{keywordBlue}\bfseries,
    stringstyle=\color{stringOrange},
    basicstyle=\ttfamily\small,
    breaklines=true,
    captionpos=b,
    keepspaces=true,
    numbers=left,
    numbersep=8pt,
    numberstyle=\tiny\color{gray},
    showstringspaces=false,
    tabsize=2,
    frame=single,
    rulecolor=\color{gray!40},
    xleftmargin=15pt,
}
\lstdefinestyle{propertiesStyle}{
    backgroundcolor=\color{lightGray},
    basicstyle=\ttfamily\small,
    breaklines=true,
    frame=single,
    rulecolor=\color{gray!40},
    numbers=left,
    numberstyle=\tiny\color{gray},
    numbersep=8pt,
    xleftmargin=15pt,
    commentstyle=\color{commentGreen},
    morecomment=[l]{\#},
}
\lstset{style=javaStyle}

% ── Tableaux ───────────────────────────────────────────────────────────────────
\usepackage{booktabs}
\usepackage{tabularx}
\usepackage{array}
\usepackage{longtable}

% ── Boites & titres ────────────────────────────────────────────────────────────
\usepackage{mdframed}
\usepackage{tcolorbox}
\tcbuselibrary{skins,breakable}
\usepackage{titlesec}

\titleformat{\chapter}[block]
  {\normalfont\huge\bfseries\color{ensetBlue}}
  {\thechapter.}{12pt}{}
\titlespacing*{\chapter}{0pt}{10pt}{20pt}

\titleformat{\section}
  {\normalfont\Large\bfseries\color{ensetBlue}}
  {\thesection}{1em}{}
\titleformat{\subsection}
  {\normalfont\large\bfseries}
  {\thesubsection}{1em}{}

% ── Diagrammes ─────────────────────────────────────────────────────────────────
\usepackage{tikz}
\usetikzlibrary{shapes.geometric, arrows.meta, positioning, fit, backgrounds}

% ── Icones statut ─────────────────────────────────────────────────────────────
\usepackage{amssymb}
\newcommand{\OK}{\textcolor{green!60!black}{$\checkmark$}}
\newcommand{\WARN}{\textcolor{orange}{$\triangle$}}
\newcommand{\ERR}{\textcolor{red}{$\times$}}

% ── En-tete / pied de page ────────────────────────────────────────────────────
\usepackage{fancyhdr}
\pagestyle{fancy}
\fancyhf{}
\fancyhead[L]{\small\textcolor{gray}{Rapport Technique -- Examen Rattrapage AD ENSET 2026}}
\fancyhead[R]{\small\textcolor{gray}{EL GHOUALI Zakariae}}
\fancyfoot[C]{\thepage}
\renewcommand{\headrulewidth}{0.4pt}

% ==============================================================================
\begin{document}

% ── Page de garde ─────────────────────────────────────────────────────────────
\begin{titlepage}
\begin{center}
    \vspace*{1cm}

    {\large\textbf{Ecole Nationale Superieure de l'Enseignement Technique}}\\[0.3cm]
    {\large ENSET -- Mohammedia}\\[0.3cm]
    {\large Filiere : Genie Logiciel / Architecture Distribuee}

    \vspace{1.5cm}
    \rule{\linewidth}{1.5pt}\\[0.4cm]
    {\Huge\bfseries\color{ensetBlue} Rapport Technique}\\[0.3cm]
    {\LARGE Examen de Rattrapage}\\[0.2cm]
    {\Large Architecture Distribuee -- JEE}\\[0.3cm]
    \rule{\linewidth}{1.5pt}

    \vspace{1.5cm}

    \begin{tcolorbox}[colback=ensetBlue!8, colframe=ensetBlue,
                      title=Contexte du projet, fonttitle=\bfseries]
        Application de gestion de credits bancaires developpee avec\\
        \textbf{Spring Boot 3.x}, \textbf{Spring Security (JWT/RSA)},\\
        \textbf{Spring Data JPA}, \textbf{H2} et \textbf{OpenAPI/Swagger}.
    \end{tcolorbox}

    \vspace{2cm}

    \begin{tabular}{rl}
        \textbf{Etudiant :}   & EL GHOUALI Zakariae \\[4pt]
        \textbf{Email :}      & zakariaelgouali@gmail.com \\[4pt]
        \textbf{Module :}     & Architecture Distribuee / JEE \\[4pt]
        \textbf{Session :}    & Rattrapage 2025--2026 \\[4pt]
        \textbf{Date :}       & \today \\
    \end{tabular}

    \vfill
    {\small Rapport technique -- Examen de rattrapage Architecture Distribuee ENSET 2026}
\end{center}
\end{titlepage}

% ── Table des matieres ────────────────────────────────────────────────────────
\tableofcontents
\newpage

% ==============================================================================
\chapter{Introduction et Presentation du Sujet}

\section{Contexte}

Cet examen de rattrapage porte sur le developpement d'une application web
backend de \textbf{gestion de credits bancaires}. Le projet met en oeuvre les
principes d'une architecture distribuee JEE en utilisant l'ecosysteme Spring
Boot.

L'objectif est de concevoir et implementer une API REST securisee permettant de :
\begin{itemize}
    \item Gerer des \textbf{clients} d'une banque et leurs credits associes ;
    \item Gerer trois types de \textbf{credits} : personnel, immobilier et
          professionnel ;
    \item Enregistrer et consulter les \textbf{remboursements} lies a chaque
          credit ;
    \item Securiser l'ensemble des endpoints via \textbf{JWT} signe avec une
          paire de cles \textbf{RSA}.
\end{itemize}

\section{Environnement Technique}

\begin{table}[h!]
\centering
\begin{tabularx}{\linewidth}{l X}
\toprule
\textbf{Technologie}         & \textbf{Role / Version} \\
\midrule
Java 21                      & Langage de programmation principal \\
Spring Boot 3.5.x            & Framework applicatif \\
Spring Data JPA              & Couche de persistance (ORM) \\
H2 Database                  & Base de donnees en memoire (developpement) \\
Spring Security              & Authentification et autorisation \\
OAuth2 Resource Server       & Validation des tokens JWT \\
Nimbus JOSE + JWT            & Encodage/decodage JWT (RSA 2048) \\
Lombok                       & Reduction du code boilerplate \\
SpringDoc OpenAPI 2.5        & Documentation Swagger UI \\
Maven                        & Outil de build et gestion des dependances \\
\bottomrule
\end{tabularx}
\caption{Stack technique du projet}
\end{table}

% ==============================================================================
\chapter{Architecture Globale}

\section{Vue d'ensemble en couches}

L'application suit une \textbf{architecture en couches} classique des
applications Spring Boot :

\begin{center}
\begin{tikzpicture}[
    layer/.style={draw=ensetBlue, fill=ensetBlue!10, rounded corners=4pt,
                  minimum width=10cm, minimum height=1cm, font=\bfseries},
    arrow/.style={-Stealth, thick, color=ensetBlue}
]
\node[layer, fill=orange!15, draw=orange] (sec)
    at (0, 1.6) {Couche Securite (Spring Security + JWT/RSA)};
\node[layer] (web)
    at (0, 0)   {Couche Web (REST Controllers)};
\node[layer] (service)
    at (0,-1.6) {Couche Service (Logique metier)};
\node[layer] (repo)
    at (0,-3.2) {Couche Repository (Spring Data JPA)};
\node[layer] (db)
    at (0,-4.8) {Base de donnees H2 (in-memory)};

\draw[arrow] (sec)     -- (web);
\draw[arrow] (web)     -- (service);
\draw[arrow] (service) -- (repo);
\draw[arrow] (repo)    -- (db);
\end{tikzpicture}
\end{center}

\section{Structure des packages}

\begin{verbatim}
ma.enset.elghouali.zakariae.exam_jee/
|-- ExamJeeApplication.java        (point d'entree + donnees de test)
|-- entities/
|   |-- Client.java
|   |-- Credit.java                (classe mere, SINGLE_TABLE)
|   |-- CreditPersonnel.java
|   |-- CreditImmobilier.java
|   |-- CreditProfessionnel.java
|   `-- Remboursement.java
|-- enums/
|   |-- StatutCredit.java
|   |-- TypeBien.java
|   `-- TypeRemboursement.java
|-- repositories/
|   |-- ClientRepository.java
|   |-- CreditRepository.java
|   `-- RemboursementRepository.java
|-- dtos/
|   |-- ClientDTO.java
|   |-- CreditDTO.java
|   `-- RemboursementDTO.java
|-- mappers/
|   `-- ExamMapper.java
|-- services/
|   |-- CreditService.java         (interface)
|   `-- CreditServiceImpl.java
|-- web/
|   |-- CreditRestController.java
|   `-- AuthController.java
`-- security/
    |-- SecurityConfig.java
    |-- RsaKeysConfig.java
    `-- PasswordEncoderConfig.java
\end{verbatim}

% ==============================================================================
\chapter{Modele de Donnees}

\section{Entites JPA}

\subsection{Client}

\begin{lstlisting}[language=Java, caption={Entite Client}]
@Entity
@Data @NoArgsConstructor @AllArgsConstructor @Builder
public class Client {
    @Id @GeneratedValue(strategy = GenerationType.IDENTITY)
    private int id;
    private String nom;
    private String email;

    @OneToMany(mappedBy = "client", fetch = FetchType.LAZY)
    private List<Credit> credits;
}
\end{lstlisting}

\subsection{Credit -- Classe mere avec heritage SINGLE\_TABLE}

\begin{lstlisting}[language=Java, caption={Entite Credit -- heritage SINGLE\_TABLE}]
@Entity
@Inheritance(strategy = InheritanceType.SINGLE_TABLE)
@DiscriminatorColumn(name = "TYPE_CREDIT", length = 15)
@Data @NoArgsConstructor @AllArgsConstructor
public class Credit {
    @Id @GeneratedValue(strategy = GenerationType.IDENTITY)
    private Long id;
    private LocalDate dateDemande;
    @Enumerated(EnumType.STRING)
    private StatutCredit statut;
    private LocalDate dateAcceptation;
    private Double montant;
    private Integer dureeRemboursement;
    private Double tauxInteret;
    @ManyToOne
    private Client client;
    @OneToMany(mappedBy = "credit", fetch = FetchType.LAZY)
    private List<Remboursement> remboursements;
}
\end{lstlisting}

\subsection{Sous-types de credits}

\begin{table}[h!]
\centering
\begin{tabularx}{\linewidth}{l l X}
\toprule
\textbf{Classe} & \textbf{Discriminateur} & \textbf{Attributs specifiques} \\
\midrule
\texttt{CreditPersonnel}
    & \texttt{PERSONNEL}
    & \texttt{motif : String} \\
\texttt{CreditImmobilier}
    & \texttt{IMMOBILIER}
    & \texttt{typeBien : TypeBien}
      (APPARTEMENT, MAISON, LOCAL\_COMMERCIAL) \\
\texttt{CreditProfessionnel}
    & \texttt{PROFESSIONNEL}
    & \texttt{motif : String}, \texttt{raisonSociale : String} \\
\bottomrule
\end{tabularx}
\caption{Sous-types de credits}
\end{table}

\subsection{Remboursement}

\begin{lstlisting}[language=Java, caption={Entite Remboursement}]
@Entity
@Data @NoArgsConstructor @AllArgsConstructor @Builder
public class Remboursement {
    @Id @GeneratedValue(strategy = GenerationType.IDENTITY)
    private Long id;
    private LocalDate date;
    private Double montant;
    @Enumerated(EnumType.STRING)
    private TypeRemboursement type;   // MENSUALITE | REMBOURSEMENT_ANTICIPE
    @ManyToOne
    private Credit credit;
}
\end{lstlisting}

\section{Enumerations du domaine}

\begin{table}[h!]
\centering
\begin{tabularx}{\linewidth}{l X}
\toprule
\textbf{Enum} & \textbf{Valeurs} \\
\midrule
\texttt{StatutCredit}      & EN\_COURS, ACCEPTE, REJETE \\
\texttt{TypeBien}          & APPARTEMENT, MAISON, LOCAL\_COMMERCIAL \\
\texttt{TypeRemboursement} & MENSUALITE, REMBOURSEMENT\_ANTICIPE \\
\bottomrule
\end{tabularx}
\caption{Enumerations du domaine}
\end{table}

\section{Strategie d'heritage et justification}

La strategie \textbf{SINGLE\_TABLE} est choisie : toutes les sous-classes de
\texttt{Credit} sont stockees dans une unique table avec une colonne discriminante
\texttt{TYPE\_CREDIT}.

\begin{itemize}
    \item \textbf{Avantage} : une seule jointure lors des requetes, performances
          elevees pour ce type de hierarchie peu profonde.
    \item \textbf{Inconvenient} : les colonnes specifiques a un sous-type sont
          nullable pour les autres sous-types.
    \item \textbf{Conclusion} : choix pertinent pour ce contexte d'examen
          (nombre limite de sous-types, hierarchie simple).
\end{itemize}

% ==============================================================================
\chapter{Couche de Persistance}

\section{Repositories Spring Data JPA}

Trois interfaces \texttt{JpaRepository} beneficient automatiquement des
operations CRUD, de la pagination et du tri :

\begin{itemize}
    \item \texttt{ClientRepository extends JpaRepository<Client, Integer>}
    \item \texttt{CreditRepository extends JpaRepository<Credit, Long>}
    \item \texttt{RemboursementRepository extends JpaRepository<Remboursement, Long>}
\end{itemize}

\section{Configuration de la base de donnees}

\begin{lstlisting}[style=propertiesStyle, caption={application.properties (version corrigee)}]
spring.application.name=exam-jee

spring.datasource.url=jdbc:h2:mem:creditdb
spring.datasource.driverClassName=org.h2.Driver
spring.datasource.username=sa
spring.datasource.password=

spring.jpa.database-platform=org.hibernate.dialect.H2Dialect
spring.jpa.hibernate.ddl-auto=update
spring.jpa.show-sql=true

spring.h2.console.enabled=true
spring.h2.console.path=/h2-console

rsa.private-key=classpath:certs/private.pem
rsa.public-key=classpath:certs/public.pem
\end{lstlisting}

\begin{tcolorbox}[colback=red!8, colframe=red!70,
                  title={\ERR\ Anomalies detectees dans application.properties}]
Le fichier original contient trois fautes de frappe bloquantes :
\begin{enumerate}
    \item \texttt{spring.datasourc\textbf{ce}.url}
          $\rightarrow$ double ``c'' ; doit etre \texttt{spring.datasource.url}
    \item \texttt{spring.datasource.password\textbf{s}}
          $\rightarrow$ ``s'' superflu ; doit etre \texttt{spring.datasource.password}
    \item \texttt{spring.\textbf{dpa}.show-sql}
          $\rightarrow$ inversion ; doit etre \texttt{spring.jpa.show-sql}
\end{enumerate}
Ces erreurs empechent le demarrage correct de l'application.
\end{tcolorbox}

\section{Initialisation des donnees de test}

Le \texttt{CommandLineRunner} dans \texttt{ExamJeeApplication} peuple la base
au demarrage avec :

\begin{itemize}
    \item 2 clients : \textit{Hassan} (hassan@gmail.com) et
          \textit{Amina} (amina@gmail.com) ;
    \item 1 credit personnel lie a Hassan :
          50\,000 DH, 48 mois, taux 4.5\%, statut \texttt{ACCEPTE},
          motif ``Achat Voiture'' ;
    \item 1 credit immobilier lie a Amina :
          1\,200\,000 DH, 240 mois, taux 4.2\%, statut \texttt{EN\_COURS},
          type bien \texttt{APPARTEMENT} ;
    \item 2 remboursements mensuels (1\,200 DH chacun) lies au credit personnel.
\end{itemize}

% ==============================================================================
\chapter{Couche Service et Mapping}

\section{Interface CreditService}

\begin{lstlisting}[language=Java, caption={Interface CreditService}]
public interface CreditService {
    List<ClientDTO>  getAllClients();
    List<CreditDTO>  getCreditsByClient(Long clientId);
    CreditDTO        getCreditById(Long creditId);
    RemboursementDTO addRemboursement(RemboursementDTO dto);
}
\end{lstlisting}

\section{Implementation CreditServiceImpl}

\begin{lstlisting}[language=Java, caption={Extrait de CreditServiceImpl}]
@Service @Transactional @AllArgsConstructor
public class CreditServiceImpl implements CreditService {

    private ClientRepository clientRepository;
    private CreditRepository creditRepository;
    private RemboursementRepository remboursementRepository;
    private ExamMapper mapper;

    @Override
    public List<CreditDTO> getCreditsByClient(Long clientId) {
        Client client = clientRepository.findById(clientId)
            .orElseThrow(() -> new RuntimeException("Client introuvable"));
        return client.getCredits().stream()
            .map(mapper::fromCredit)
            .collect(Collectors.toList());
    }

    @Override
    public RemboursementDTO addRemboursement(RemboursementDTO dto) {
        Credit credit = creditRepository.findById(dto.getCreditId())
            .orElseThrow(() -> new RuntimeException("Credit introuvable"));
        Remboursement r = Remboursement.builder()
            .date(dto.getDate()).montant(dto.getMontant())
            .type(dto.getType()).credit(credit).build();
        return mapper.fromRemboursement(remboursementRepository.save(r));
    }
}
\end{lstlisting}

\section{Mapper ExamMapper}

Le mapper utilise \texttt{BeanUtils.copyProperties} pour les proprietes
communes, puis determine le type du credit par polymorphisme
(\texttt{instanceof}) afin de renseigner \texttt{typeCredit} dans le DTO.

\begin{lstlisting}[language=Java, caption={Methode fromCredit du mapper}]
public CreditDTO fromCredit(Credit credit) {
    CreditDTO dto = new CreditDTO();
    BeanUtils.copyProperties(credit, dto);
    if (credit instanceof CreditPersonnel)
        dto.setTypeCredit("PERSONNEL");
    else if (credit instanceof CreditImmobilier)
        dto.setTypeCredit("IMMOBILIER");
    else if (credit instanceof CreditProfessionnel)
        dto.setTypeCredit("PROFESSIONNEL");
    return dto;
}
\end{lstlisting}

\begin{tcolorbox}[colback=orange!10, colframe=orange,
                  title={\WARN\ Anomalie detectee dans ExamMapper}]
La valeur retournee pour \texttt{CreditProfessionnel} etait \texttt{"PROFESSION"}
au lieu de \texttt{"PROFESSIONNEL"} dans le code source. Cela cree une
incoherence avec la valeur du discriminateur JPA
(\texttt{@DiscriminatorValue("PROFESSIONNEL")}). La version corrigee est
presentee ci-dessus.
\end{tcolorbox}

% ==============================================================================
\chapter{API REST}

\section{Tableau des endpoints}

\begin{longtable}{l l l p{5.5cm}}
\toprule
\textbf{Methode} & \textbf{URL}
    & \textbf{Auth} & \textbf{Description} \\
\midrule
\endhead
POST & \texttt{/auth/login}                & Non & Authentification et generation du JWT \\
GET  & \texttt{/api/clients}               & JWT & Liste de tous les clients \\
GET  & \texttt{/api/clients/\{id\}/credits}& JWT & Credits d'un client \\
GET  & \texttt{/api/credits/\{id\}}        & JWT & Detail d'un credit par id \\
POST & \texttt{/api/remboursements}         & JWT & Enregistrement d'un remboursement \\
\bottomrule
\caption{Endpoints REST de l'application}
\end{longtable}

\section{Controleur principal}

\begin{lstlisting}[language=Java, caption={CreditRestController}]
@RestController
@RequestMapping("/api")
@CrossOrigin("*")
@AllArgsConstructor
public class CreditRestController {

    private CreditService creditService;

    @GetMapping("/clients")
    public List<ClientDTO> clients() {
        return creditService.getAllClients();
    }

    @GetMapping("/clients/{id}/credits")
    public List<CreditDTO> creditsByClient(
            @PathVariable(name = "id") Long clientId) {
        return creditService.getCreditsByClient(clientId);
    }

    @GetMapping("/credits/{id}")
    public CreditDTO getCredit(@PathVariable Long id) {
        return creditService.getCreditById(id);
    }

    @PostMapping("/remboursements")
    public RemboursementDTO saveRemboursement(
            @RequestBody RemboursementDTO dto) {
        return creditService.addRemboursement(dto);
    }
}
\end{lstlisting}

\section{Controleur d'authentification}

\begin{lstlisting}[language=Java, caption={AuthController -- generation du JWT}]
@PostMapping("/auth/login")
public Map<String, String> login(String username, String password) {
    // 1. Authentification via AuthenticationManager
    Authentication auth = authenticationManager.authenticate(
        new UsernamePasswordAuthenticationToken(username, password));

    // 2. Construction du scope (roles)
    String scope = auth.getAuthorities().stream()
        .map(GrantedAuthority::getAuthority)
        .collect(Collectors.joining(" "));

    // 3. Claims JWT : validite 30 minutes
    JwtClaimsSet claims = JwtClaimsSet.builder()
        .issuedAt(Instant.now())
        .expiresAt(Instant.now().plus(30, ChronoUnit.MINUTES))
        .subject(username)
        .claim("scope", scope)
        .build();

    // 4. Signature RSA et retour
    String jwt = jwtEncoder
        .encode(JwtEncoderParameters.from(claims)).getTokenValue();
    return Map.of("access-token", jwt);
}
\end{lstlisting}

% ==============================================================================
\chapter{Securite}

\section{Architecture de securite}

\begin{center}
\begin{tikzpicture}[
    box/.style={draw=ensetBlue, fill=ensetBlue!8, rounded corners=4pt,
                minimum width=3.2cm, minimum height=1cm,
                align=center, font=\small},
    arr/.style={-Stealth, thick, color=ensetBlue}
]
\node[box] (client)  at (0, 0)    {Client\\(app / curl)};
\node[box] (login)   at (5, 1.2)  {POST /auth/login\\username + password};
\node[box] (tok)     at (10, 1.2) {JWT signe\\RSA 2048};
\node[box] (api)     at (5, -1.2) {GET /api/**\\Bearer token};
\node[box] (filter)  at (10,-1.2) {JwtDecoder\\verification RSA};

\draw[arr] (client.east) |- (login.west);
\draw[arr] (login) -- node[above,font=\tiny]{access-token} (tok);
\draw[arr] (tok.west) -- (login.east);
\draw[arr] (client.east) |- (api.west);
\draw[arr] (api) -- (filter);
\draw[arr] (filter.east) -- ++(1.2,0) node[right,font=\tiny]{200 OK / 401};
\end{tikzpicture}
\end{center}

\section{Configuration Spring Security}

\begin{lstlisting}[language=Java, caption={SecurityConfig -- chaine de filtres}]
@Bean
public SecurityFilterChain filterChain(HttpSecurity http) throws Exception {
    return http
        .csrf(csrf -> csrf.disable())
        .authorizeHttpRequests(auth -> auth
            .requestMatchers(
                "/auth/**", "/swagger-ui.html", "/swagger-ui/**",
                "/v3/api-docs/**", "/h2-console/**").permitAll())
        .authorizeHttpRequests(auth -> auth.anyRequest().authenticated())
        .sessionManagement(sess ->
            sess.sessionCreationPolicy(SessionCreationPolicy.STATELESS))
        .oauth2ResourceServer(oauth2 ->
            oauth2.jwt(Customizer.withDefaults()))
        .build();
}
\end{lstlisting}

\section{Utilisateurs en memoire}

\begin{table}[h!]
\centering
\begin{tabularx}{\linewidth}{l l l X}
\toprule
\textbf{Utilisateur} & \textbf{Mot de passe} & \textbf{Role(s)} & \textbf{Droits} \\
\midrule
\texttt{client1}  & 1234 & ROLE\_CLIENT              & Consultation \\
\texttt{employe1} & 1234 & ROLE\_EMPLOYE             & Gestion remboursements \\
\texttt{admin}    & 1234 & ROLE\_ADMIN, ROLE\_EMPLOYE & Acces complet \\
\bottomrule
\end{tabularx}
\caption{Utilisateurs InMemoryUserDetailsManager}
\end{table}

\section{Cles RSA}

Les cles RSA sont generees hors de l'application et placees dans
\texttt{src/main/resources/certs/} :
\begin{itemize}
    \item \texttt{private.pem} -- cle privee RSA 2048 bits (signature des tokens) ;
    \item \texttt{public.pem}  -- cle publique RSA (verification des tokens).
\end{itemize}
La liaison est assuree par le record \texttt{RsaKeysConfig} active via
\texttt{@EnableConfigurationProperties(RsaKeysConfig.class)}.

% ==============================================================================
\chapter{Documentation API -- Swagger / OpenAPI}

L'integration de \textbf{SpringDoc OpenAPI 2.5} expose automatiquement :
\begin{itemize}
    \item \texttt{/swagger-ui.html} -- interface graphique interactive ;
    \item \texttt{/v3/api-docs}     -- specification OpenAPI JSON.
\end{itemize}
Ces deux URLs sont explicitement autorisees sans authentification dans
\texttt{SecurityConfig}, ce qui permet de les consulter librement lors du
developpement.

% ==============================================================================
\chapter{Analyse Qualite et Points d'Amelioration}

\section{Bilan des anomalies detectees}

\begin{longtable}{l p{4.8cm} p{5cm}}
\toprule
\textbf{Severite} & \textbf{Probleme} & \textbf{Correction recommandee} \\
\midrule
\endhead

\ERR\ CRITIQUE
    & 3 fautes de frappe bloquantes dans \texttt{application.properties}
      (\texttt{datasourcce}, \texttt{passwords}, \texttt{dpa})
    & Corriger : \texttt{datasource.url}, \texttt{datasource.password},
      \texttt{jpa.show-sql} \\[6pt]

\WARN\ MOYEN
    & \texttt{ExamMapper} retourne \texttt{"PROFESSION"} au lieu de
      \texttt{"PROFESSIONNEL"} pour \texttt{CreditProfessionnel}
    & Corriger la chaine en \texttt{"PROFESSIONNEL"} \\[6pt]

\WARN\ MOYEN
    & \texttt{@CrossOrigin("*")} trop permissif (CORS ouvert)
    & Restreindre aux origines connues en production \\[6pt]

\WARN\ MOYEN
    & Exceptions generiques \texttt{RuntimeException} dans la couche service
    & Creer \texttt{ClientNotFoundException}, \texttt{CreditNotFoundException} \\[6pt]

\WARN\ MOYEN
    & Aucune pagination sur \texttt{/api/clients}
    & Ajouter \texttt{Pageable} et retourner \texttt{Page<ClientDTO>} \\[6pt]

\WARN\ MINEUR
    & Les roles definis dans \texttt{SecurityConfig} ne sont pas utilises
      dans les endpoints via \texttt{@PreAuthorize}
    & Ajouter \texttt{@PreAuthorize} sur chaque endpoint \\[6pt]

\OK\ INFO
    & \texttt{FetchType.LAZY} sur toutes les associations
    & Bonne pratique, evite le N+1 implicite \\[6pt]

\OK\ INFO
    & SINGLE\_TABLE adapte au nombre de sous-types
    & Choix pertinent pour ce contexte \\[6pt]

\OK\ INFO
    & JWT RSA 2048 + duree 30 min
    & Securite asymetrique correcte \\

\bottomrule
\caption{Bilan qualite du projet}
\end{longtable}

\section{Ameliorations suggeries}

\subsection{Gestion centralisee des erreurs}

\begin{lstlisting}[language=Java, caption={GlobalExceptionHandler a ajouter}]
@RestControllerAdvice
public class GlobalExceptionHandler {

    @ExceptionHandler(ClientNotFoundException.class)
    @ResponseStatus(HttpStatus.NOT_FOUND)
    public Map<String, String> handleClientNotFound(
            ClientNotFoundException ex) {
        return Map.of("error", ex.getMessage());
    }

    @ExceptionHandler(CreditNotFoundException.class)
    @ResponseStatus(HttpStatus.NOT_FOUND)
    public Map<String, String> handleCreditNotFound(
            CreditNotFoundException ex) {
        return Map.of("error", ex.getMessage());
    }
}
\end{lstlisting}

\subsection{Securite au niveau des methodes}

\begin{lstlisting}[language=Java, caption={Controle par role via @PreAuthorize}]
@GetMapping("/clients")
@PreAuthorize("hasRole('EMPLOYE') or hasRole('ADMIN')")
public List<ClientDTO> clients() { ... }

@PostMapping("/remboursements")
@PreAuthorize("hasRole('EMPLOYE') or hasRole('ADMIN')")
public RemboursementDTO saveRemboursement(@RequestBody RemboursementDTO dto)
    { ... }
\end{lstlisting}

\subsection{Validation des entrees}

\begin{lstlisting}[language=Java, caption={Validation sur RemboursementDTO}]
public class RemboursementDTO {
    @NotNull
    private LocalDate date;
    @Positive
    private Double montant;
    @NotNull
    private TypeRemboursement type;
    @NotNull
    private Long creditId;
}
\end{lstlisting}

% ==============================================================================
\chapter{Conclusion}

\section{Recapitulatif des realisations}

Ce projet d'examen de rattrapage implemente une application bancaire de gestion
de credits couvrant l'ensemble des fonctionnalites demandees :

\begin{itemize}
    \item[\OK] Modelisation JPA complete avec heritage \texttt{SINGLE\_TABLE}
               (3 types de credits) ;
    \item[\OK] Couche de persistance Spring Data JPA avec H2 ;
    \item[\OK] Couche service transactionnelle et mapper manuel ;
    \item[\OK] API REST documentee avec SpringDoc/OpenAPI ;
    \item[\OK] Authentification JWT asymetrique RSA 2048 via Spring Security ;
    \item[\OK] Separation propre en couches
               (Entity $\to$ Repository $\to$ Service $\to$ Controller) ;
    \item[\WARN] Fautes de configuration dans \texttt{application.properties}
                 a corriger ;
    \item[\WARN] Robustesse a ameliorer : exceptions metier, pagination,
                 validation, \texttt{@PreAuthorize} ;
    \item[\ERR] Tests unitaires et d'integration absents.
\end{itemize}

\section{Competences demontrees}

\begin{table}[h!]
\centering
\begin{tabularx}{\linewidth}{X l}
\toprule
\textbf{Competence} & \textbf{Statut} \\
\midrule
Conception et modelisation JPA / ORM        & \OK\ Maitrisee \\
Heritage JPA (SINGLE\_TABLE)               & \OK\ Maitrisee \\
Spring Data JPA (repositories)             & \OK\ Maitrisee \\
Developpement REST avec Spring MVC         & \OK\ Maitrisee \\
Securisation JWT / OAuth2 Resource Server  & \OK\ Maitrisee \\
Cryptographie asymetrique (RSA)            & \OK\ Maitrisee \\
Pattern DTO + Mapper                       & \OK\ Maitrisee \\
Documentation API (OpenAPI / Swagger)      & \OK\ Maitrisee \\
Gestion d'exceptions centralisee           & \WARN\ Partielle \\
Validation des entrees (\texttt{@Valid})   & \WARN\ Absente \\
Tests unitaires / integration              & \ERR\ Absents \\
\bottomrule
\end{tabularx}
\caption{Bilan des competences JEE}
\end{table}

\vspace{1cm}

\begin{tcolorbox}[colback=ensetBlue!8, colframe=ensetBlue,
                  title={Appreciation globale}]
Le projet demontre une bonne maitrise de l'ecosysteme Spring Boot et des
principes d'architecture distribuee JEE. La structure en couches est claire,
la securisation par JWT/RSA est pertinente et bien configuree. Les anomalies
relevees (typos de configuration, valeur enum incorrecte, absence de validation
et de tests) sont des points de perfectionnement courants a ce niveau
d'etudes.
\end{tcolorbox}

\end{document}
